\documentclass[11pt]{article}
\usepackage[margin=1in]{geometry}
\usepackage{amsmath,amssymb,amsthm}
\usepackage{booktabs}

\newtheorem{theorem}{Theorem}
\newtheorem{corollary}{Corollary}
\newtheorem{lemma}{Lemma}
\newcommand{\F}{\mathbb{F}}
\newcommand{\EPL}{\operatorname{EPL}}

\title{V75: Symbolic Prefix Counting on Bounded-Width Gate Decompositions}
\author{NC0\_k-Avoid Laboratory}
\date{1 August 2026}

\begin{document}
\maketitle

\begin{abstract}
We compile the exact weighted affine-residual dynamic program of Laboratory V74
into a monotone arithmetic circuit for a paired-variable output generating
polynomial.  On a supplied gate branch decomposition of boundary width $b$,
the circuit has $O(mA(b)^2)$ arithmetic operations.  Maintaining gate values
under paired-variable updates gives a depth-sensitive avoided-output search in
$O(A(b)^2(m+\sum_i\operatorname{depth}_T(i))\operatorname{poly}(n,m))$ time.
A supplied logarithmic-height decomposition therefore yields an
$O(m\log m\,A(b)^2\operatorname{poly}(n,m))$ algorithm.  Caterpillar trees show
that this does not improve the arbitrary supplied-tree worst case.  No
unrestricted Range Avoidance or P-versus-NP consequence is claimed.
\end{abstract}

\section{Model}
Let $C:\{0,1\}^n\to\{0,1\}^m$ be represented by Boolean output gates of
fan-in at most three.  For every output coordinate $i$, introduce formal
variables $z_{i,0}=u_i$ and $z_{i,1}=v_i$, and define
\[
 P_C(u,v)=\sum_{x\in\{0,1\}^n}\prod_{i=1}^m z_{i,C_i(x)}.
\]
Every selected gate fiber is represented by an exact pairwise-disjoint union
of affine cells, as in V74.  Let $T$ be a rooted binary tree whose leaves are
the output gates.  At node $t$, let $B_t$ be the variables used by gates both
inside and outside the leaf set of $t$.  Assume $|B_t|\le b$ at every node.
Let $A(b)$ denote the number of nonempty affine subspaces of $\F_2^b$.

\section{Exact polynomial}
\begin{theorem}[Exact paired generating polynomial]
For every $y\in\{0,1\}^m$, the coefficient of
$\prod_i z_{i,y_i}$ in $P_C$ is exactly $|C^{-1}(y)|$.  If $p$ is an output
prefix, setting the selected variable of each fixed coordinate to one, its
opposite to zero, and both variables of every unfixed coordinate to one gives
the exact prefix count $N(p)$.
\end{theorem}

\begin{proof}
Each input $x$ contributes exactly one monomial, namely the monomial encoding
$C(x)$.  Therefore the multiplicity of the monomial encoding $y$ equals the
number of inputs mapping to $y$.  Under the prefix substitution, a monomial
survives with value one exactly when its encoded output extends $p$.
\end{proof}

\section{Residual compilation}
For each node $t$ and each projected affine residual $R$ on $B_t$, construct an
arithmetic expression $Q_t(R)$.  At a gate leaf $i$, a cell $F$ in output fiber
$c$ contributes
\[
 2^{\dim(F)-\dim(\operatorname{proj}_{B_i}F)}z_{i,c}
\]
to the expression indexed by its projection.  At an internal node with
children $u,v$, compatible residuals $R_u,R_v$ intersect to $J$ and contribute
\[
 Q_u(R_u)Q_v(R_v)
 2^{\dim(J)-\dim(\operatorname{proj}_{B_t}J)}
\]
to the expression indexed by $\operatorname{proj}_{B_t}J$.

\begin{theorem}[Monotone residual circuit]
The root empty-boundary expression, multiplied by the free-unused-input factor,
is exactly $P_C$.  The construction uses only nonnegative integer constants,
addition, and multiplication.  It creates
\[
 S=O(mA(b)^2)
\]
arithmetic operations, apart from polynomial-time affine intersection and
projection.
\end{theorem}

\begin{proof}
Induct on $T$.  For a node $t$, residual $R$, subtree output word $y_t$, and
boundary assignment satisfying $R$, the coefficient of the monomial selected
by $y_t$ in $Q_t(R)$ equals the number of compatible internal assignments per
such boundary assignment.  At a leaf, affine projection has uniform fiber size
$2^{\dim(F)-\dim(\operatorname{proj}F)}$, and disjoint cells prevent double
counting.  At a join, residual intersection enforces consistency of the two
children, multiplication combines independent internal choices, and the
projection factor counts eliminated assignments uniformly.  At the root the
boundary is empty, so the coefficients are global preimage counts.

There are at most $A(b)$ residuals per node and at most $A(b)^2$ child pairs per
join.  A binary tree with $m$ leaves has $O(m)$ nodes, and every pair creates
only constantly many arithmetic operations.
\end{proof}

\section{Incremental evaluation}
Store every arithmetic gate value and reverse dependency edges.  Let $D_T(i)$
be the number of arithmetic operation nodes reachable from either paired
variable of output coordinate $i$.

\begin{theorem}[Depth-sensitive incremental evaluation]
An avoided output for $m>n$ can be constructed using
\[
 O\!\left(S+\sum_{i=1}^m D_T(i)\right)
\]
arithmetic reevaluations.  Moreover,
\[
 D_T(i)=O\!\left(A(b)^2\operatorname{depth}_T(i)\right),
\]
so the total running time, including affine operations, is
\[
 O\!\left(A(b)^2\left(m+\sum_i\operatorname{depth}_T(i)\right)
 \operatorname{poly}(n,m)\right).
\]
\end{theorem}

\begin{proof}
Initially all paired variables equal one and the root value is $2^n$.  At
coordinate $i$, set $v_i$ to zero and propagate only through its reverse
dependency cone; the resulting root is $N(p0)$.  Since the two exact gate
fibers partition the parent prefix, $N(p1)=N(p)-N(p0)$.  Choose a child whose
count is below its remaining completion capacity.  If bit one is selected,
restore $v_i$ and set $u_i$ to zero in one batch.  Each coordinate therefore
causes at most two dependency-cone reevaluations.  The standard prefix
pigeonhole invariant gives a zero-preimage output after $m$ coordinates.

Only operations created at leaf $i$ and at ancestors of that leaf can depend
on $u_i$ or $v_i$.  Each ancestor contributes at most $O(A(b)^2)$ operations,
which proves the cone bound.
\end{proof}

\begin{corollary}[Balanced supplied decomposition corollary]
If $T$ has height $h$, then the running time is
\[
 O(mhA(b)^2\operatorname{poly}(n,m)).
\]
In particular, a supplied height-$O(\log m)$ decomposition gives
\[
 O(m\log m\,A(b)^2\operatorname{poly}(n,m)).
\]
\end{corollary}

\section{Depth obstruction}
For a left-deep caterpillar on $m\ge2$ leaves,
\[
 \EPL(T)=\sum_i\operatorname{depth}_T(i)=\frac{m^2+m-2}{2}.
\]
Thus local dependency-cone reuse alone retains quadratic worst-case work on an
arbitrary supplied decomposition.  The laboratory proves no balancing theorem
with controlled width inflation.

\section{Finite verification}
The primary verifier checks 4,096 exhaustive two-input, three-output circuits,
32,768 coefficients, 61,440 prefixes, and 4,096 incrementally constructed
avoided outputs.  A seeded varied-support ternary suite checks 48 circuits on
balanced and caterpillar trees, totaling 6,144 coefficient and 12,192 prefix
comparisons.  An independent verifier reconstructs direct truth-table
semantics without importing the arithmetic-circuit or affine-residual code.

\section{Scope}
This result assumes a supplied bounded-width gate decomposition.  It does not
prove automatic width-preserving balancing, unrestricted
$\mathrm{NC}^0_3$-Avoid, a superpolynomial lower bound, a size-preserving
simulation to a standard proof or branching model, novelty, peer review, or any
consequence for P versus NP.

\end{document}
